\documentclass[dvipdfmx]{jarticle}
\hoffset=-25mm \voffset=-25mm
\setlength{\textwidth}{16cm}
\setlength{\textheight}{23cm}
\usepackage{url}
\usepackage{tikz}
\usetikzlibrary{calc,shapes}
\newcounter{coffee}
\newcounter{bean}
\begin{document}
\begin{flushright}
  \texttt{2026/10/16}
\end{flushright}
\begin{center}{\bf\Large
    2026年度　コンパイラ演習課題(2)
  }\end{center}
\begin{itemize}
  \item TACTの課題の添付ファイルとしてPDFファイルで提出すること．
  \item 手書きレポートをPDFとして提出方法は講義webサイトを
        参照\footnote{\url{https://sites.google.com/sqlab.jp/26compiler/how-to-make-report}}すること．
  \item 提出期限: 10月16日(金) 23:55(JST)
\end{itemize}
\begin{list}{[\arabic{coffee}]}{\usecounter{coffee}}
  \item つぎの算術式を逆ポーランド記法(解析木のポストフィックス記法)で表せ．
        \begin{list}{(\alph{bean})}{\usecounter{bean}}
          \item \texttt{a*b+c}
          \item \texttt{a*b-c*(d+e)}
        \end{list}
  \item CF文法 $G_1=\langle\{E,S\},\{i,t.e,a,b\},P_1,S\}\rangle$について以下の問いに答えよ．

        \[P_1=\left\{\begin{array}{rcl}
            S & \rightarrow & i\,E\,t\,S,       \\
            S & \rightarrow & i\,E\,t\,S\,e\,S, \\
            S & \rightarrow & a,                \\
            E & \rightarrow & b
          \end{array}\right\}\]
        \begin{list}{(\alph{bean})}{\usecounter{bean}}
          \item $G_1$に左括りだしを適用した文法の生成規則を示せ．
          \item 上記の文法を$G'_1$とすると，$G'_1$が曖昧であることを示せ．
          \item $G_1$, $G'_1$に対して下向き構文解析表を作成せよ．
        \end{list}
  \item 以下のCF文法$G_2$に対して以下の問いに答えよ．
        \[G_2=\langle\{S,B,D,E,F\},\{u,v,w,x,y,z\},P_2,S\rangle\]

        \[P_2=\left\{\begin{array}{ll}
            S\rightarrow uBDz , & E\rightarrow y          , \\
            B\rightarrow Bv   , & E\rightarrow\varepsilon , \\
            B\rightarrow w    , & F\rightarrow x          , \\
            D\rightarrow EF   , & F\rightarrow\varepsilon
          \end{array}\right\}\]

        \begin{list}{(\alph{bean})}{\usecounter{bean}}
          \item $\alpha\in PO_{P_2}\cup\{S,B,D,E,F\}$に対する$\mathsf{First}_1(\alpha)$を求めよ．
          \item $A\in\{S,B,D,E,F\}$に対する$\mathsf{Follow}_1(A)$を求めよ．
          \item 下向き構文解析表を作成し，LL(1)文法ではないことを示せ．
          \item 左再帰性を除去して等価なLL(1)文法$G'_2$を構成し，下向き構文解析表を作成せよ．
          \item $G'_2$による$uvvyxz$に対する最左導出を示し，構文解析木を示せ．ここで，$G'_2$の生成規則に番号$1,2,\cdots$を付けて，規則$n$による導出を${\Rightarrow}_{n}$と書くこと．
        \end{list}
\end{list}
\end{document}
